\scanpage{19}
\mspage{1}

\diagram{\begin{tikzcd}
Y\arrow[r,hook,"i"]\arrow[d,"g"']&X\arrow[d,"f"]&U\arrow[l,hook]\arrow[d,"\wr"]\\
Y'\arrow[r,hook,"i'"']&X'&U'\arrow[l,hook]
\end{tikzcd}}

Soient $X'$ dans $(\mathrm{Cat})$

$Y'$ sous-objet fermé, $U'$ compl. ouvert

\margin{\unsure{Même}}

\[
f:X\longrightarrow X',\quad Y=f^{-1}(Y'),\qquad U=f^{-1}(U')
\]

Supposons que $U\xrightarrow{\sim}U'$

alors \deleted{$f$}\inserted{le carré} est-il néc. cocartésien ? Non bien sûr, comme on voit en prenant $\unsure{X'=U'|X=U}$, $f$ l'inclusion, alors $Y=\varnothing$.

\uline{Mais} supposons que $f$ soit cofibré, le carré est-il cocartésien ? En d'autres termes, a-t-on pour $u'\in U'$, $y'\in Y'$
\[
\mathrm{Hom}_{X'}(u',y')\xleftarrow[\sim]{?}
\underbrace{\underset{y\;\mathrm{ds}\;Y_{/y'}}{\varinjlim}\mathrm{Hom}(u,y)}_{\underset{Y_{y'}}{\varinjlim}\mathrm{Hom}(u,Y_{y'})}\;?
\]
\margin{(où $u\in\mathrm{Ob}\,U$ au dessus de $u'$) ?}

\diagram{\begin{tikzcd}[column sep=large,row sep=large]
u&\deleted{\unsure{y}}\arrow[d,dashed,dash]&y\arrow[d,dashed]\\
u'\arrow[r,dashed,dash]&\deleted{\unsure{f(y)}}\arrow[r]&y'
\end{tikzcd}}

Oui, car \deleted{$\mathrm{Hom}(u,$}

(i) Surjectif et

(ii) Pour $\unsure{\xi|f}:u\longrightarrow y'$ fixé \ldots

Nous nous intéressons donc à la situation des diagrammes du type précédent, avec $f$ \uline{cocartésien} / Peut se décrire par un foncteur
\[
X'\longrightarrow(\mathrm{Cat}),
\]
\deleted{dans} dont la restriction à $U'$ soit le foncteur $c^t$ de valeur $\Delta_0=e_{(\mathrm{Cat})}$,

i.e.

(i) $Y'\xrightarrow{\Phi}\mathrm{Cat}$ (définissant \deleted{$\unsure{f}$} $Y\longrightarrow Y'$)

(ii) $\mathrm{Hom}(u',y')\longrightarrow\mathrm{Ob}\,\Phi(y')$ bifonctoriel en $u',y'$ ds $U'\times Y'$
% archive foliation: 18
